\chapter{Introduction}
\label{chap:introduction}

Music festivals and large outdoor events have become a massive part of modern culture. Every summer, millions of people gather in open fields, parks, and temporary venues to enjoy live performances and food. From Untold in Romania to Tomorrowland in Belgium, these events keep growing in size and complexity. But behind all the excitement, there is a very practical problem that organizers have to deal with: how do you let tens of thousands of people buy food, drinks, and merchandise in a place where the internet barely works?\\

The traditional approach was straightforward: people brought cash, stood in line, and paid at the counter. It was slow, it was messy with always having change and there was always the risk of theft or losing your wallet in the crowd. Over the past decade, many festivals have moved towards cashless payment systems, typically based on RFID wristbands or prepaid cards. The idea is simple: you load money onto your wristband before or during the event, and then you just tap it at a vendor terminal to pay. Industry data from cashless technology providers suggests that RFID systems reduce transaction times by up to 80 percent compared to cash and that festival organizers see average revenue increases of around 22 percent after going cashless \cite{Tappit2024}. These figures are consistent with the academic evidence: a recent meta-analysis across 71 studies found a statistically significant, small positive effect of cashless payments on consumer spending ($g = 0.135$, 95\% CI [0.068; 0.203]), driven primarily by the reduced psychological friction of not physically handing over money \cite{Schomburgk2024}.

However, these systems come with a critical weakness: most of them depend on a reliable internet connection. And if there is one thing that festivals are known for, it is terrible connectivity. When tens of thousands of people gather in the same area, all using their phones simultaneously, the cellular network simply cannot handle the load. Shafiq et al.~\cite{Shafiq2013} conducted the first large-scale characterization of cellular network performance during crowded events using real-world traces from a tier-1 US network. Their findings are striking: pre-connection failures increased by 100 to 5,000 times compared to routine days, driven by users exhausting the limited bandwidth of the signaling channel when too many devices attempt to acquire radio resources simultaneously. This degradation extended as far as 10 miles around the event venue as attendees arrived and departed. Even for users who managed to connect, voice call failures increased by 7 to 30 times, packet, and round-trip times increased by up to 7 times. The performance spikes were particularly severe during intermissions, when large numbers of attendees simultaneously attempted to use their devices. Notably, the study found that even with portable base stations and free Wi-Fi offloading already deployed, these degradation levels persisted, highlighting a fundamental capacity limitation rather than a simple configuration problem.


This is where things get problematic for payment systems. If the payment terminal at a food stall needs to contact a remote server to authorize a transaction, and the network is down or painfully slow, the transaction fails. The attendee cannot buy food. The vendor loses a sale. Multiply this by tens of thousands of transactions per day, and you can see why this is a serious issue. The most famous example of this going wrong happened at the Download Festival in the UK in 2015, the first major UK festival to go completely cashless. The RFID system experienced technical failures that left attendees unable to purchase food or drinks, with many taking to social media to describe the system as useless. Although organisers claimed only around one percent of attendees were affected, the backlash was severe enough that the festival abandoned the cashless-only approach the following year \cite{eFestivals2016}.

The core question this thesis tries to answer is: can we build a payment system that works reliably at festivals, and that does not fall apart when the network has a bad moment? Our approach is to design a system that operates online by default, backed by a dedicated private Wi-Fi network deployed at the event venue. This private network serves as the primary connectivity layer, keeping vendor terminals connected to the cloud backend for real-time transaction processing. But because no network is perfectly reliable, especially one assembled in a field for a few days, the system also includes a modular offline extension that kicks in when connectivity is temporarily lost. The offline module is not the core of the architecture, it is a safety net. Certain operations, most notably wallet top-ups, are inherently online because they involve moving real money through external providers and cannot be handled locally.

\section{Problem Statement}
\label{sec:problem-statement}

Festival environments present a unique combination of challenges that are not commonly addressed together in existing payment system literature. The starting point is connectivity. As discussed above, cellular networks collapse under the load of mass gatherings, and many festivals take place in rural or temporary venues where wired infrastructure simply does not exist. Our approach addresses this by deploying a dedicated private Wi-Fi network on site, which serves as the primary connectivity layer between vendor terminals and the cloud backend. Under normal conditions, the system operates fully online: transactions are authorized in real time against a central database, wallet balances are always up to date, and operations like top-ups that depend on external payment providers (Stripe, in our case) work without any issues. The private network effectively sidesteps the public cellular congestion that causes so many problems for existing solutions.

However, even a well-designed private network is not immune to failure. Equipment can malfunction, interference from the sheer density of wireless devices in a small area can cause temporary dead zones, and multi-day events put sustained stress on infrastructure that was assembled in a field a week earlier. When the network goes down, even for a few minutes during a peak period, the system still needs to process payments. This is where the offline extension module comes in. Rather than treating offline operation as the primary design goal, we treat it as a resilience layer: a modular component that activates when connectivity is temporarily lost and ensures that the festival economy does not grind to a halt. The key word here is temporarily, because prolonged offline operation introduces complications that are best avoided, as we will discuss below.

The second challenge is that these are closed ecosystems. Unlike a regular retail environment where a customer walks into a shop and pays with a bank card, a festival creates its own little economy. Attendees load credits, spend them at authorized vendors, and ideally get a refund for whatever they did not spend. This is what the payment industry calls a closed-loop system: the money circulates within a defined perimeter, controlled by a single operator, rather than flowing through the global banking network. The closed-loop nature is actually an advantage in our case, because it means we control the entire transaction flow and do not depend on external authorization networks that would be unreachable during a connectivity outage.

The third challenge is scale and speed. A major festival might have hundreds of vendor points, processing thousands of transactions per hour, especially during peak times like meal breaks or between performances. The payment system must be fast enough that it does not create bottlenecks. Attendees are not going to wait two minutes for a payment to go through when they just want a sandwich. In the online mode, the cloud backend handles this load centrally. In the offline mode, each terminal processes transactions locally, which is actually faster per transaction but introduces the next challenge.

The fourth challenge is consistency during offline periods. When terminals lose connectivity and start processing transactions independently, they make local decisions about whether a payment is valid without being able to check with the central server or with each other. What happens when two terminals approve transactions that together exceed a user's balance? This is the classic double-spend problem, and it is one of the hardest issues to solve in distributed systems. For a festival payment system, refusing to process transactions because we cannot guarantee consistency is not an acceptable answer, because that means people cannot buy food. So during these offline windows, we choose availability and accept that some inconsistencies may occur that need to be resolved once connectivity is restored.

This leads to the fifth challenge: reconciliation. When terminals reconnect after an offline period, all the transactions they processed independently need to be merged back into the central system state. The synchronization protocol must ensure that no transactions are lost and that conflicts, such as a balance that went negative due to concurrent offline spending, are resolved in a deterministic way. Because our system is online-first, these reconciliation windows are expected to be relatively short and infrequent compared to a system that operates offline by default. But they still need to be handled correctly, because even a brief network outage during a peak period can generate hundreds of transactions that need to be synchronized.

It is also worth noting that not all operations can work offline. Wallet top-ups, for instance, require real-time communication with Stripe (or any other On Ramp provider) and the cloud backend, because actual money is moving from the attendee's bank account into the system. You cannot process that locally. This is an intentional design boundary: the offline module handles payments (debiting the wallet), but funding the wallet (crediting it) remains an online-only operation.

\section{Objectives}
\label{sec:objectives}

The main objective of this thesis is to design and implement a minimum viable product (MVP) for a cloud-based, closed-loop payment system tailored to festival and event environments, with a modular offline extension that provides resilience during temporary connectivity loss. The concept of a minimum viable product, as defined by Ries~\cite{Ries2011}, refers to a version of a product that allows the team to collect the maximum amount of validated learning with the least effort. In our case, this means building enough of the system to demonstrate that the core architectural ideas work, without necessarily implementing every feature a production deployment would need. Throughout this work, we explore different design architectures and patterns, discussing their benefits and trade-offs in the context of festival environments, before settling on the approach that best fits the constraints of the problem.

Specifically, the objectives of this work are the following:

\textbf{O1:} Design a system architecture that operates online by default through a dedicated private Wi-Fi network deployed at the festival site, with all core payment operations processed in real time against a cloud backend. The architecture should also include a modular offline extension that activates automatically when connectivity is temporarily lost, allowing vendor terminals to continue processing payments locally.

\textbf{O2:} Implement a wallet-based payment service capable of processing debit and credit operations between consumer and vendor accounts. Wallet top-ups are handled via Stripe and require an active connection to the cloud backend. Debit transactions (payments at vendor terminals) should work both online, where they are authorized centrally, and offline, where they are processed locally until connectivity is restored.

\textbf{O3:} Explore and evaluate different strategies for keeping the on-site and cloud components in sync, particularly after periods of offline operation. The chosen approach should ensure that transactions processed during connectivity loss are eventually reflected in the central system and that conflicts are handled in a consistent way.

\textbf{O4:} Implement a core service that manages vendors, products, stock, users, and events, providing the administrative backbone of the system.

\textbf{O5:} Develop client applications for the three main user roles: a Consumer App for attendees to top up their wallet and make payments via QR or NFC, a Vendor App for merchants to manage products and accept payments, and an Admin Dashboard for event organizers to oversee the entire operation.

\textbf{O6:} Evaluate the proposed architecture against the specific requirements of festival environments, discussing the trade-offs between online and offline modes, the limitations of the current implementation, and the design decisions that shaped the system.

\section{Related Work}
\label{sec:related-work}

The idea of building payment systems that work without constant connectivity is not new, but most of the existing work comes from domains that are quite different from festivals. Transit card systems like Suica, Oyster, and Octopus have been doing offline transactions for decades, storing balances directly on the card chip so that gates can open without waiting for a server response. These systems are well studied and provide useful lessons about transaction speed and security pitfalls, but they operate in a controlled infrastructure with fixed terminals and predictable usage patterns, which is not exactly what a festival looks like.

On the software architecture side, the local-first movement, particularly the work by Kleppmann et al.~\cite{Kleppmann2019}, lays out principles for applications that treat the local device as the primary copy and synchronize with the cloud when they can. Concepts like Conflict-Free Replicated Data Types offer formal guarantees about convergence in distributed systems, and event sourcing provides a natural model for logging transactions that happened offline and replaying them later. These ideas inform our design, though they do not map one-to-one onto payment systems, since money has the annoying property of being a conserved quantity.

On the commercial side, providers like Tappit, Intellitix, Weezevent, and Glownet have deployed cashless systems at hundreds of festivals, but their architectures are proprietary and essentially undocumented in academic literature. The failures we do know about, such as the Download Festival 2015 crash or the Danish Nets outage in 2024, highlight the consequences of getting this wrong, but they do not tell us much about how to get it right.

In short, the building blocks exist across transit systems, distributed systems theory, and commercial festival deployments, but nobody has put them together into a documented, open architecture designed for the specific constraints of festival environments. That is the gap this thesis tries to fill. A more detailed treatment of each of these areas follows in Chapter~\ref{chap:background}.
